% ---- pg-415 （书页 403） ----

\mvbold{评注}\quad 可以很简单地构造具有好的直观性质的因子得分. 将在一个因子上有高载荷(例如绝对值大于 0.40)的变量归为一组. 可以通过把组内变量的(标准化)观测值根据载荷的符号求和得到因子得分. 通过这些简单的因子得分来代替标准化数据可以使数据简化. 简单因子得分与用最小二乘或回归方法得到的因子得分之间通常有高度相关性.

\noindent\textbf{例 9.13}\quad (由因子分析分组创建简单综合得分)

例 9.4 中股票价格数据主成分因子分析的载荷估计为
\[
\tilde{\mathbf{L}}=
\begin{bmatrix}
0.732&-0.437\\
0.831&-0.280\\
0.726&-0.374\\
0.605&0.694\\
0.563&0.719
\end{bmatrix}
\qquad\text{与}\qquad
\tilde{\mathbf{L}}^*=\tilde{\mathbf{L}}\mathbf{T}=
\begin{bmatrix}
0.852&0.030\\
0.851&0.214\\
0.813&0.079\\
0.133&0.911\\
0.084&0.909
\end{bmatrix}
\]
对每个因子, 把 $\tilde{\mathbf{L}}$ 中具有最大绝对值的载荷看成同等大小, 忽略小的载荷. 这样我们创建综合变量为线性组合
\[
\hat{f}_1=x_1+x_2+x_3+x_4+x_5
\]
\[
\hat{f}_2=x_4+x_5-x_1
\]
在实践中, 我们可以标准化这些新变量.

如果我们不是用 $\tilde{\mathbf{L}}$, 而是用方差最大旋转载荷矩阵 $\tilde{\mathbf{L}}^*$, 则简单因子得分为
\[
\hat{f}_1=x_1+x_2+x_3
\]
\[
\hat{f}_2=x_4+x_5
\]
确定高载荷和可忽略载荷是相当主观的判断. 最好选择与主题内容有关的线性组合. \hfill$\blacksquare$

虽然在因子分析中常常假定变量为多元正态的, 但是对大量的变量来说, 证明这个假定是很困难的. 像我们在第 4 章指出的那样, 边际变换可以帮助我们. 类似的, 因子得分可以是也可以不是正态分布的, 双变量因子得分散布图可以产生各种各样非椭圆形状的图形. 在运用因子得分作其他分析之前, 应对其图形作检查. 它们可以揭示出远离中心的值和(可能的)非正态性的程度.

\section{因子分析的展望和建议}
\phantomsection\label{sec:9.6}

在因子分析研究中, 必须作出许多决策. 大概最重要的决策是选择公共因子数 $m$. 虽然已知 $m$ 时对模型充分性有大样本检验, 但这只是对近似正态分布的数据适用. 而且, 若变量数和观测数都大, 检验将会很有把握地拒绝极小 $m$ 的模型, 而正是此时因子分析可提供一个有用的近似. 最为常见的是, $m$ 最后的选定基于下述考虑的综合: (1)所解释的样本方差的比例, (2)题材知识, 以及 (3)结论的“合理程度”.

解法和旋转类型的选择并非决定性的决策, 事实上, 大多数满意的因子分析都试用了不止一种旋转并且所有结果实质上确认同一因子结构.

在现阶段, 因子分析仍然是一种艺术, 并且尚没有一种单一的策略可以获得完美的成功. 我们建议和详细阐述一种合情合理的选择:

\noindent 1.\quad 执行主成分因子分析. 这个方法对第一次考察数据尤其合适(不要求 $\mathbf{R}$ 或 $\mathbf{S}$ 是非奇异的).
